\documentclass[11pt]{article}

\usepackage[margin=1in]{geometry}
\usepackage[T1]{fontenc}
\usepackage{lmodern}
\usepackage{microtype}
\usepackage{amsmath,amssymb,amsthm,mathtools}
\usepackage[dvipsnames]{xcolor}
\usepackage[most]{tcolorbox}
\usepackage{tikz}
\usepackage{enumitem}
\usepackage{tocloft}
\usepackage[hidelinks]{hyperref}

\definecolor{courseblue}{HTML}{1F4E79}
\definecolor{lightblue}{HTML}{EAF2F8}
\definecolor{lightgreen}{HTML}{EAF7EA}
\definecolor{lightpink}{HTML}{FCE8EF}
\definecolor{lightyellow}{HTML}{FFF8D6}
\definecolor{warningyellow}{HTML}{D6A000}
\definecolor{exercisebg}{HTML}{D6E6F5}
\definecolor{exerciseframe}{HTML}{1F4E79}
\definecolor{examplebg}{HTML}{F4FAFE}
\definecolor{exampleframe}{HTML}{85C1E9}
\definecolor{darkgray}{HTML}{333333}

\setlength{\parindent}{0pt}
\setlength{\parskip}{0.65em}
\setlist[itemize]{topsep=0.25em,itemsep=0.15em}
\renewcommand{\cftsecleader}{\cftdotfill{\cftdotsep}}
\renewcommand{\cftsubsecleader}{\cftdotfill{\cftdotsep}}

\theoremstyle{plain}
\newtheorem{theorem}{Theorem}[section]
\newtheorem{proposition}[theorem]{Proposition}
\newtheorem{lemma}[theorem]{Lemma}
\newtheorem{corollary}[theorem]{Corollary}
\theoremstyle{definition}
\newtheorem{definition}[theorem]{Definition}
\newtheorem{example}[theorem]{Example}
\newtheorem{exercise}[theorem]{Exercise}
\theoremstyle{remark}
\newtheorem*{remark}{Remark}
\newtheorem*{fact}{Fact}
\newtheorem*{warning}{Warning}

\tcolorboxenvironment{definition}{
  enhanced,
  colback=lightgreen,
  colframe=ForestGreen,
  boxrule=0.7pt,
  arc=2pt,
  left=6pt,
  right=6pt,
  top=5pt,
  bottom=5pt
}

\tcolorboxenvironment{warning}{
  enhanced,
  colback=lightyellow,
  colframe=warningyellow,
  boxrule=0.7pt,
  arc=2pt,
  left=6pt,
  right=6pt,
  top=5pt,
  bottom=5pt
}

\tcolorboxenvironment{exercise}{
  enhanced,
  colback=exercisebg,
  colframe=exerciseframe,
  boxrule=1pt,
  arc=2pt,
  left=6pt,
  right=6pt,
  top=5pt,
  bottom=5pt
}

\tcolorboxenvironment{example}{
  enhanced,
  colback=examplebg,
  colframe=exampleframe,
  boxrule=0.7pt,
  arc=2pt,
  left=6pt,
  right=6pt,
  top=5pt,
  bottom=5pt
}

\numberwithin{equation}{section}

\newcommand{\be}{\begin{eqnarray}}
\newcommand{\ee}{\end{eqnarray}}
\newcommand{\C}{\mathbb{C}}
\newcommand{\R}{\mathbb{R}}
\newcommand{\N}{\mathbb{N}}

\DeclareMathOperator{\Hol}{\mathcal{H}\mathrm{ol}}
\DeclareMathOperator{\Real}{Re}
\DeclareMathOperator{\Imag}{Im}
\newcommand{\A}{\mathcal{A}}
\newcommand{\CT}{\C[[T]]}
\newcommand{\CLT}{\C((T))}
\tcbset{breakable}
\begin{document}

\begin{center}
  \colorbox{lightblue}{%
    \begin{minipage}{0.94\textwidth}
      \vspace{0.55em}
      \begin{tabular*}{\textwidth}{@{\extracolsep{\fill}}lr}
        {\color{courseblue}\bfseries\LARGE Lecture 7} &
        {\color{courseblue}\bfseries\LARGE MATH 411}\\[0.35em]
        \textbf{Fall 2026} & \textbf{Date: Sep 18, 2026}\\
        & \textbf{Instructor: Manish Patnaik}
      \end{tabular*}
      \vspace{0.45em}
    \end{minipage}}
\end{center}

\setcounter{tocdepth}{1}
\tableofcontents
\vspace{0.45em}

% Handwritten page 1.
\section{Recollections from last time}
\begin{itemize}
\item Convergence of a sequence $(z_n)$, $z_n\in\C$, and of series
$\sum a_n$, $\sum|a_n|$.
\item For $S\subseteq\C$ and $f:S\to\C$ bounded, the sup norm
\be
 \|f\|_S=\sup_{z\in S}|f(z)|,
\ee

\end{itemize}

Given a sequence $(f_n)$ of functions on $S$, what does $f_n\to f$ mean?

\begin{definition}
	Let $f_n,f:S\to\C$.
	\begin{itemize}
		\item We say $f_n$ \emph{converges pointwise} to $f$ if $f_n(z)\to f(z)$
		for all $z\in S$.
		\item We say $f_n$ \emph{converges uniformly} to $f$ if, given
		$\epsilon>0$, there exists $n_0$ such that if $n\ge n_0$, then
		\be
		|f_n(z)-f(z)|<\epsilon
		\qquad\text{for every }z\in S.
		\ee
		Equivalently, for every $\epsilon>0$, there exists $n_0$ such that
		$\|f_n-f\|_S<\epsilon$ whenever $n\ge n_0$.
	\end{itemize}
\end{definition}
%
%\begin{definition}
%	The sequence $\{f_n\}$ is \emph{uniformly Cauchy} if, given $\epsilon>0$,
%	there is $n_0$ such that $|f_n(z)-f_m(z)|<\epsilon$ for every $z\in S$
%	whenever $m,n\ge n_0$. Equivalently, $\|f_n-f_m\|_S<\epsilon$.
%\end{definition}

\begin{remark}
	The difference between the two notions of convergence is the order of the
	quantifiers: for pointwise convergence the index $n_0$ may depend on the
	point $z\in S$, whereas for uniform convergence a single $n_0$ must work
	for every $z\in S$ at once.
\end{remark}
%
%\begin{itemize}
%\item \emph{Pointwise convergence}: for any $z\in S$,
%$f_n(z)\to f(z)$, i.e.\ $|f(z)-f_n(z)|\to0$ as $n\to\infty$.
%\item \emph{Uniform convergence}: $\|f-f_n\|_S\to0$ as $n\to\infty$.
%\end{itemize}

\section{Uniform limits of continuous functions}
\begin{theorem}
\label{thm:unif-cts}
Suppose $(f_n)$ is a sequence of continuous functions on $S$ and $f_n\to f$
uniformly. Then $f$ is also continuous.
\end{theorem}

% Handwritten page 2.
\begin{proof}
Fix $z\in S$. We want to show $f$ is continuous at $z$, i.e.
\be
 \lim_{w\to z}f(w)=f(z):
\ee
given $\epsilon>0$, there exists $\delta>0$ such that if $|w-z|<\delta$ then
$|f(z)-f(w)|<\epsilon$.

\textbf{Step 1.} Given $\epsilon/3$, by uniformity there exists $n_0$ such
that if $n\ge n_0$ then
\be
 \|f-f_n\|_S<\frac\epsilon3.
\ee

\textbf{Step 2.} Since each $f_n$ is continuous, pick some $n\ge n_0$;
continuity of $f_n$ at $z$ means there exists $\delta>0$ such that if
$|z-w|<\delta$ then
\be
 |f_n(z)-f_n(w)|<\frac\epsilon3.
\ee

\textbf{Step 3.} Now for that same $\delta$, we know that whenever
$|z-w|<\delta$,
\begin{align*}
 |f(z)-f(w)|
 &\le\underbrace{|f(z)-f_n(z)|}_{\le\ \epsilon/3\ \text{(uniformity)}}
 +\underbrace{|f_n(z)-f_n(w)|}_{<\ \epsilon/3\ \text{($f_n$ continuous)}}
 +\underbrace{|f_n(w)-f(w)|}_{\le\ \epsilon/3\ \text{(uniformity)}}\\
 &<\epsilon.
\end{align*}
Hence $f$ is continuous at $z$, and since $z\in S$ was arbitrary, $f$ is
continuous on $S$.
\end{proof}

% Handwritten page 3.
\section{Uniformly Cauchy sequences}
Another useful notion is the following:

\begin{definition}
Let $(f_n)$ be a sequence of functions on $S$. We say that $(f_n)$ is
\emph{uniformly Cauchy on $S$} if, given $\epsilon>0$, there exists $n_0$
such that if $n,m\ge n_0$, then
\be
 \|f_n-f_m\|_S<\epsilon.
\ee
\end{definition}
\begin{lemma}
\label{lem:unif-cauchy}
Let $(f_n)$ be uniformly Cauchy on $S$. Then
\begin{enumerate}
\item[(a)] $\displaystyle f(z)=\lim_{n\to\infty}f_n(z)$ exists for each
$z\in S$;
\item[(b)] $f_n\to f$ uniformly.
\end{enumerate}
\end{lemma}

\begin{proof}
\textbf{(a)} Use that, for fixed $z\in S$,
\be
 |f_n(z)-f_m(z)|\le\|f_n-f_m\|_S.
\ee
So $(f_n(z))$ is a Cauchy sequence of complex numbers, and hence
converges. Call the limit $f(z)$.

\textbf{(b)} Recall that $f_n\to f$ uniformly means: given $\epsilon>0$
there exists $n_0$ such that if $n\ge n_0$ then $\|f_n-f\|_S<\epsilon$.

Since $(f_n)$ is uniformly Cauchy, there exists $N=N(\epsilon)$ such that
\be
 \|f_n-f_m\|_S<\frac\epsilon2\qquad\text{if }n,m\ge N.
\ee
% Handwritten page 4.
Fix $n\ge N$ and $z\in S$. For every $m\ge N$ we have
\be
 |f_n(z)-f_m(z)|\le \|f_n-f_m\|_S<\frac\epsilon2.
\ee
Letting $m\to\infty$ and using part (a), we obtain
\be
 |f_n(z)-f(z)|\le\frac\epsilon2.
\ee
This holds for every $z\in S$, so
$\|f_n-f\|_S\le\epsilon/2<\epsilon$ whenever $n\ge N$. Thus $f_n\to f$
uniformly.
\end{proof}

% Handwritten page 5.
\section{Series of functions}
\subsection*{Applications}
What we are interested in is convergence of power series of the form
\be
 \sum a_nz^n,\qquad a_n\in\C,
\ee
whose partial sums
\be
 S_n=\sum_{j=0}^{n}a_jz^j
\ee
are functions on $S$; we are interested in when they converge to a function
$f$ on $S$.

\begin{definition}[Convergence of a series of functions]
Let $(f_j)_{j\ge1}$ be a sequence of functions on $S$, and put
\be
 S_d=\sum_{j=1}^{d}f_j.
\ee
We say that $\sum f_j$ converges \emph{pointwise} (respectively,
\emph{uniformly}) to $f$ if $S_d\to f$ pointwise (respectively, uniformly).
We say that $\sum f_j$ converges \emph{absolutely on $S$} if
\be
 \sum_{j=1}^{\infty}|f_j(z)|<\infty
 \qquad\text{for every }z\in S.
\ee
If the series of functions $\sum |f_j|$ converges uniformly, we say that
$\sum f_j$ converges \emph{uniformly absolutely} on $S$.
\end{definition}

% Handwritten page 6.
\begin{theorem}[Comparison test for series of functions]
\label{thm:M-test}
Let $(c_n)$ be a sequence of nonnegative real numbers such that
$\displaystyle\sum_{n=1}^{\infty}c_n$ converges. Let $(f_n)$ be a
sequence of functions on $S$ satisfying
\be
 \|f_n\|_S\le c_n,
\ee
Then $\sum f_n$ and $\sum|f_n|$ converge uniformly on $S$. In particular,
$\sum f_n$ converges absolutely on $S$. If, moreover, the $f_n$ are
continuous, then the limit is continuous.
\end{theorem}

\begin{proof}
For $m<n$ the partial sums satisfy
\be
 \|S_n-S_m\|_S=\Bigl\|\sum_{j=m+1}^{n}f_j\Bigr\|_S
 \le\sum_{j=m+1}^{n}\|f_j\|_S\le\sum_{j=m+1}^{n}c_j,
\ee
which is small for $m,n$ large because $\sum c_n$ converges (its partial
sums are Cauchy). So $(S_n)$ is uniformly Cauchy, and
Lemma~\ref{lem:unif-cauchy} gives a limit $f$ with $S_n\to f$ uniformly.
The same estimate applied to $\sum|f_j|$, using
$\||f_j|\|_S=\|f_j\|_S\le c_j$, gives uniform convergence of
$\sum|f_j|$. Hence $\sum f_j$ converges absolutely at every point of $S$.
Continuity of the limit follows from Theorem~\ref{thm:unif-cts}.
\end{proof}

\begin{theorem}
\label{thm:power-series-disc}
Let $(a_n)$ be a sequence of complex numbers and let $r>0$. If
\be
 \sum|a_n|r^n\quad\text{converges,}
\ee
then $\sum a_nz^n$ converges absolutely and uniformly on
\be
 D(0,r)=\{z\in\C:|z|<r\}
\ee
to a continuous function.
\end{theorem}

% Handwritten page 7.
\begin{proof}
Put $f_n(z)=a_nz^n$ on $D(0,r)$. Then
\be
 \|f_n\|_{D(0,r)}\le|a_n|r^n.
\ee
Use the comparison test (Theorem~\ref{thm:M-test}) with $c_n=|a_n|r^n$:
it follows that $\sum a_nz^n$ converges uniformly and absolutely on
$D(0,r)$, and the limit is continuous since each $f_n$ is.
\end{proof}

\begin{example}
Consider
\be
 \sum_{n=0}^{\infty}\frac{z^n}{n!},
 \qquad\text{so } a_nz^n=\frac{z^n}{n!},\quad a_n=\frac1{n!}.
\ee
\textbf{Claim.} The exponential series $\sum a_nr^n$ converges for every
$r>0$.

\emph{Proof:} compare to a geometric series. Fix a positive integer
$N\ge2r$. For
$n>N$,
\be
 \frac{r^n}{n!}=\frac{r^N}{N!}\cdot\frac{r^{n-N}}{(N+1)(N+2)\cdots n}
 \le\frac{r^N}{N!}\left(\frac rN\right)^{n-N}
 \le\frac{r^N}{N!}\left(\frac12\right)^{n-N},
\ee
and $\sum(1/2)^{n-N}$ converges. So by the comparison test for series of
complex numbers (Lecture 6) the series converges.

Therefore, by Theorem~\ref{thm:power-series-disc},
\be
 \boxed{\ e^z:=\sum_{n=0}^{\infty}\frac{z^n}{n!}\ }
\ee
exists for all $z\in\C$, and defines a continuous function on $\C$.
\end{example}

% Handwritten page 8.
\section{Radius of convergence}
Let $f(T)\in\CT$, say
\be
 f(T)=\sum a_nT^n.
\ee

\begin{theorem}[Radius of convergence]
\label{thm:radius}
Given any $f(T)\in\CT$, there exists $r$ with $0\le r\le\infty$ such that
\begin{itemize}
\item $\sum a_nz^n$ converges absolutely for $|z|<r$, and uniformly on
$D(0,s)$ for every $0<s<r$;
\item $\sum a_nz^n$ diverges for all $|z|>r$.
\end{itemize}
\end{theorem}
%
%\begin{proof}
%Assume $f(T)$ converges for some $z\in\C$ but not for all $z\in\C$ (the
%two extreme cases give $r=0$ and $r=\infty$). Let
%\be
% r=\text{l.u.b.\ of the set of all }s>0\text{ such that }
% \sum|a_n|s^n\text{ converges.}
%\ee
%We know $0<r<\infty$.
%
%If $z\in\C$ with $|z|>r$, then $\sum|a_n|\,|z|^n$ diverges, by the
%definition of the least upper bound. So $\sum a_nz^n$ cannot converge
%absolutely for this $|z|$.
%
%
%In fact the series does not just diverge absolutely, but it diverges genuinely as well.  Indeed, suppose $\sum a_nw^n$ converged for some $w$. Then its terms are
%bounded, say $|a_nw^n|\le M$ for all $n$. For any $z$ with $|z|<|w|$,
%\be
%\sum|a_n|\,|z|^n
%=\sum|a_nw^n|\left(\frac{|z|}{|w|}\right)^n
%\le M\sum\left(\frac{|z|}{|w|}\right)^n<\infty,
%\ee
%a convergent geometric series. Taking $z=s$ real shows that
%$\sum|a_n|s^n$ converges for every $s<|w|$, so $r\ge s$ for all such $s$,
%i.e.\ $|w|\le r$. Contrapositively, if $|w|>r$ then $\sum a_nw^n$ diverges.
%
%
%If $|z|<r$, then by definition of the least upper bound there is $s$ with
%$|z|<s\le r$ such that $\sum|a_n|s^n$ converges. But then, by
%Theorem~\ref{thm:power-series-disc},
%\be
% \sum a_nz^n\quad\text{converges uniformly and absolutely on }D(0,s),
%\ee
%and in particular at our $z$. Since every $z$ with $|z|<r$ lies in such a
%disc, $\sum a_nz^n$ converges absolutely on $D(0,r)$ and uniformly on each
%$D(0,s)$ with $s<r$.
%\end{proof}
%




\begin{remark}
Uniform convergence on all of $D(0,r)$ can fail, which is why the statement
is made on the discs $D(0,s)$, $s<r$. For instance $\sum z^n$ has $r=1$, and
the convergence is uniform on each $D(0,s)$ with $s<1$, but not on $D(0,1)$.
\end{remark}

\end{document}
